% !PRINTED 340
\weberpage{340}
%<<CONT

3. \gesperrt{Wenn $\varDelta$ positiv ist, soll $\sqrt{\varDelta}$ gleichfalls positiv
sein, und wenn $\varDelta$ negativ ist, soll $-i\sqrt{\varDelta}$ positiv, oder
$\sqrt{\varDelta}$ positiv imagin\"ar sein.}

Aus (1), (2) ergibt sich noch:

4. \gesperrt{Die Grundzahl eines quadratischen K\"orpers ist
immer $\equiv 0$ oder $\equiv 1$ nach dem Modul 4 und hat, au\ss er 4,
keinen quadratischen Teiler. Sie ist also Diskriminante
und zwar Stammdiskriminante.}

\textbf{\S\ 91. \gesperrt{Ideale und Formen in quadratischen K\"orpern.}}

Ist $\mathfrak{a}$ ein Ideal des quadratischen K\"orpers $\Omega$, $\alpha_1, \alpha_2$ eine
positive Basis, und
\weq{1}{\lambda=\alpha_1 t_1+\alpha_2 t_2}
eine Basisform von $\mathfrak{a}$, so ist
\weq{2}{N(\lambda)=(\alpha_1 t_1+\alpha_2 t_2)(\alpha_1' t_1+\alpha_2' t_2)=N(\mathfrak{a})\,T,}
und $N(\mathfrak{a})$ ist der gr\"o\ss te gemeinschaftliche Teiler der drei ganzen
rationalen Zahlen
\[
\alpha_1\alpha_1', \quad \alpha_1\alpha_2'+\alpha_2\alpha_1', \quad \alpha_2\alpha_2'.
\]
Setzen wir also
\weq{3}{\begin{aligned}
\alpha_1\alpha_1' &= a\,N(\mathfrak{a}),\\
\alpha_1\alpha_2'+\alpha_2\alpha_1' &= b\,N(\mathfrak{a}),\\
\alpha_2\alpha_2' &= c\,N(\mathfrak{a}),
\end{aligned}}
so ist
\weq{4}{T=a t_1^2+b t_1 t_2+c t_2^2}
die zur Basis $\alpha_1, \alpha_2$ geh\"orige Form. Wir bezeichnen sie, wenn
es auf die Bezeichnung der Variablen nicht ankommt, mit
\weq{5}{\varphi=(a, b, c).}
Es ist eine primitive quadratische Form, deren Diskriminante
\weq{6}{\varDelta=b^2-4ac}
gleich der Grundzahl des K\"orpers $\Omega$ ist (\S\ 87). Um die Basis
$\alpha_1, \alpha_2$ nach Bd.\,II, \S\ 163 zu bilden, suche man zun\"achst die kleinste
durch $\mathfrak{a}$ teilbare nat\"urliche Zahl $a_{11}$ und hierauf die kleinste
nat\"urliche Zahl $a_{12}$, f\"ur die sich eine rationale, der Bedingung
\weq{7}{a_{12}+a_{22}\theta\equiv 0\pmod{\mathfrak{a}}}
gen\"ugende Zahl $a_{12}$ bestimmen % [sic] In beiden Vorkommen steht a_{12}.
 l\"a\ss t. Da die Kongruenz (7) f\"ur
$a_{22}=a_{11}$ befriedigt werden kann, n\"amlich durch $a_{12}=0$, so
%CONT>>
